../cictikz/src/cictikz/data/tex/cictikz_preamble.tex

% The architecture of a low power radio: a shared matching network, a
% receiver (LNA, mixer, anti-alias filters, ADCs), the all-digital PLL
% that makes the local oscillator, and the transmitter.

\begin{circuitikz}[american, thick, transform shape,
  blk/.style={draw, rectangle, minimum width=1.1cm, minimum height=0.9cm}]

  \def\yI{5.3}
  \def\yQ{3.9}

  % the antenna
  \draw (0,5.0) -- (0,5.8);
  \draw (-0.4,6.4) -- (0,5.8) -- (0.4,6.4);
  \draw[-latex] (0,5.0) -- (0.7,5.0);

  % the matching network, shared by RX and TX
  \draw (0.7,0.6) rectangle (1.9,5.2);
  \node[anchor=center] at (1.3,2.9) {\small MATCH};

  % the RX front end
  \node[blk] (lna) at (2.9,5.0) {LNA};
  \draw[-latex] (1.9,5.0) -- (lna.west);

  \draw (4.0,3.4) rectangle (4.8,5.6);
  \node[anchor=center] at (4.4,4.5) {MIX};
  \draw[-latex] (lna.east) -- (4.0,5.0);

  % the two baseband branches
  \node[blk] (aafi) at (5.9,\yI) {AAF};
  \node[blk] (adci) at (7.5,\yI) {ADC};
  \node[blk] (aafq) at (5.9,\yQ) {AAF};
  \node[blk] (adcq) at (7.5,\yQ) {ADC};

  \draw[-latex] (4.8,\yI) -- (aafi.west);
  \draw[-latex] (aafi.east) -- (adci.west);
  \draw[-latex] (adci.east) -- ++(0.8,0);
  \node[anchor=west] at (8.95,\yI) {$I_{RX}$};

  \draw[-latex] (4.8,\yQ) -- (aafq.west);
  \draw[-latex] (aafq.east) -- (adcq.west);
  \draw[-latex] (adcq.east) -- ++(0.8,0);
  \node[anchor=west] at (8.95,\yQ) {$Q_{RX}$};

  % the all digital PLL that makes the LO
  \draw (4.0,2.4) rectangle (8.05,3.1);
  \node[anchor=center] at (6.0,2.75) {AD PLL};
  \draw[-latex] (9.0,2.75) -- (8.05,2.75);
  \node[anchor=west] at (9.1,2.75) {$f_{REF}$};
  \draw[latex-latex] (4.4,3.1) -- (4.4,3.4);

  % the transmitter
  \draw (2.6,0.6) rectangle (8.05,2.0);
  \node[anchor=center] at (5.3,1.3) {TX};
  \draw[latex-latex] (4.4,2.0) -- (4.4,2.4);
  \draw[-latex] (9.0,1.6) -- (8.05,1.6);
  \node[anchor=west] at (9.1,1.6) {$I_{TX}$};
  \draw[-latex] (9.0,1.0) -- (8.05,1.0);
  \node[anchor=west] at (9.1,1.0) {$Q_{TX}$};
  \draw[-latex] (2.6,1.3) -- (1.9,1.3);

\end{circuitikz}
\end{document}
